% Section V: Classical solutions, level quantization, and Wilson observables
% Purpose: Global structure of Chern-Simons theory
%
% Four-layer organization:
% 1. Classical theory: action and F_A = 0
% 2. Global structure: large gauge transformations and level quantization via π₃(G) ≅ ℤ
% 3. Quantization: phase space is moduli space of flat G-connections on Σ
% 4. Observables: Wilson loops, linking, knot data, framing
%
% Status: MERGE AND REORGANIZE scattered content
% Sources:
% - tex_docs/chern_simons_two_review_synthesis.tex (Schwartz current, level quantization)
%   Lines ~1-500: CS action, level quantization, Wilson loops - GOOD pedagogical base
% - tex_docs/dense_derivation_expansion.tex (worked calculations)
%   Lines ~200-800: Level quantization calculation, Wilson loop phases, linking - DETAILED
% - writings/brian_chern_simons_theory.tex (mathematical rigor)
%   Sections on flat connections, holonomy, moduli spaces - RIGOROUS
%
% Action needed: MERGE following the four-layer structure
% Use synthesis.tex for narrative, dense_derivation for calculations, Brian for rigor

\section{Classical Structure and Observables}

% TODO: Organize into four clear layers below

\subsection{Classical Solutions: Flat Connections}

% TODO: Equation of motion F_A = 0
% Solution space: flat connections
% Source: chern_simons_two_review_synthesis.tex, lines ~100-200
% Add: Math background on holonomy and flat connections HERE (from topology_review.tex)

% Introduce π₁(M) and holonomy here naturally:
% "Flat connections are classified by holonomy around non-contractible cycles"

\subsection{Large Gauge Transformations and Level Quantization}

% TODO: This is crucial - where π₃(G) enters
% Source:
% - chern_simons_two_review_synthesis.tex, lines ~300-450 (pedagogical)
% - dense_derivation_expansion.tex, lines ~400-600 (detailed calculation)
% - brian_chern_simons_theory.tex (rigorous proof)
%
% Include π₃(G) ≅ ℤ background HERE (from topology_review.tex)
% Show that level k must be integer
% This is the quantization condition

% FIGURE FLAG [F-level-quantization]: Diagram showing gauge transformation on S³,
% winding number from π₃(G), and requirement k ∈ ℤ

\subsection{Moduli Space of Flat Connections}

% TODO: M_flat(M,G) ≅ Hom(π₁(M),G)/G for closed 3-manifold boundaries
% Phase space structure for canonical quantization
% Source: brian_chern_simons_theory.tex (best treatment)
%
% Include symplectic geometry of moduli space HERE (from topology_review.tex)
% This prepares for Section VI quantization

\subsection{Wilson Loop Observables}

% TODO: W_R(K) = tr_R 𝒫exp(∮_K A)
% Gauge invariance
% Linking and framing
% Source:
% - chern_simons_two_review_synthesis.tex (introduction)
% - dense_derivation_expansion.tex, lines ~650-800 (explicit calculation for abelian case)
%
% Show abelian Gaussian integral calculation from dense_derivation
% Reference nonabelian case (Jones polynomial, quantum groups) - don't fully develop

% FIGURE FLAG [F-wilson-linking]: Two linked loops with Wilson lines,
% showing how CS theory computes linking number
